\documentclass[aspectratio=169]{beamer}

\usepackage[utf8]{inputenc}
\usepackage[T1]{fontenc}
\usepackage[brazil]{babel}
\usepackage{amsmath, amssymb}
\usepackage{microtype}
\usepackage{csquotes}
\usepackage{natbib}
\usepackage{tikz}
\usetikzlibrary{positioning, arrows.meta}

% Colors
\definecolor{mblack}{RGB}{20, 20, 20}
\definecolor{mprimary}{RGB}{44, 7, 53}     % dark blue
\definecolor{maccent}{RGB}{133, 138, 227}    % orange
\definecolor{mwhite}{RGB}{255, 255, 255}

% Theme
\usetheme{default}
\usecolortheme{default}
\setbeamercolor{normal text}{fg=mblack, bg=white}
\setbeamercolor{frametitle}{fg=white, bg=mprimary}
\setbeamercolor{title}{fg=white}
\setbeamercolor{author}{fg=white}
\setbeamercolor{date}{fg=white}
\setbeamercolor{itemize item}{fg=maccent}
\setbeamercolor{itemize subitem}{fg=mprimary}
\setbeamercolor{enumerate item}{fg=maccent}
\setbeamercolor{enumerate subitem}{fg=mprimary}
\setbeamercolor{block title}{fg=white, bg=mprimary}
\setbeamercolor{block body}{fg=mblack, bg=mprimary!8}
\setbeamercolor{block title alerted}{fg=white, bg=maccent}
\setbeamercolor{block body alerted}{fg=mblack, bg=maccent!10}
\setbeamercolor{alerted text}{fg=maccent}

\setbeamerfont{frametitle}{series=\bfseries, size=\large}
\setbeamerfont{title}{series=\bfseries, size=\LARGE}
\setbeamerfont{block title}{series=\bfseries}

\setbeamertemplate{navigation symbols}{}

% Footer
\setbeamertemplate{footline}{%
  \leavevmode%
  \hbox{%
    \begin{beamercolorbox}[wd=0.7\paperwidth, ht=2.25ex, dp=1ex, leftskip=0.3cm]{normal text}%
      \tiny\color{mblack!50} Filosofia da Ciência -- Lakatos e Kuhn
    \end{beamercolorbox}%
    \begin{beamercolorbox}[wd=0.3\paperwidth, ht=2.25ex, dp=1ex, rightskip=0.3cm, right]{normal text}%
      \tiny\color{mblack!50}\insertframenumber{} / \inserttotalframenumber
    \end{beamercolorbox}%
  }%
}

% Frametitle with accent rule
\setbeamertemplate{frametitle}{%
  \nointerlineskip%
  \begin{beamercolorbox}[wd=\paperwidth, ht=0.55cm, dp=0.2cm, leftskip=0.5cm]{frametitle}%
    \usebeamerfont{frametitle}\insertframetitle%
  \end{beamercolorbox}%
  \vspace{-0.1cm}%
  \textcolor{maccent}{\rule{\paperwidth}{2pt}}%
}

% Title page: three-band layout
\setbeamercolor{titlebar top}{bg=mwhite}
\setbeamercolor{titlebar main}{bg=mprimary}
\setbeamercolor{titlebar bottom}{bg=maccent}

\setbeamertemplate{title page}{%
  \vspace{0pt}%
  \begin{beamercolorbox}[wd=\paperwidth, ht=0.55cm, dp=0pt]{titlebar top}\end{beamercolorbox}%
  \begin{beamercolorbox}[wd=\paperwidth, ht=3.8cm, dp=0pt, center]{titlebar main}%
    \vspace{0.6cm}
    {\usebeamerfont{title}\color{white}\inserttitle\par}%
    \vspace{0.35cm}{\color{white}\small\insertauthor\par}%
    \vspace{0.15cm}{\color{white}\footnotesize\insertdate\par}%
    \vspace{0.25cm}
  \end{beamercolorbox}%
  \begin{beamercolorbox}[wd=\paperwidth, ht=1.3cm, dp=0pt]{titlebar bottom}\end{beamercolorbox}%
}

\setlength{\itemsep}{4pt}

% Highlight commands
\newcommand{\impt}[1]{\textcolor{maccent}{\textbf{#1}}}
\newcommand{\ntc}[1]{\textcolor{mprimary}{#1}}

\title{Teorias com estrutura: Lakatos e Kuhn}
\author{Filosofia da Ciência\footnote{Baseado em \citet{chalmers1993}, caps.~VII e VIII.}}
\date{\today}

\begin{document}

\begin{frame}[plain]\titlepage\end{frame}

\begin{frame}{Sumário}\tableofcontents\end{frame}

% =====================================================
\section{Recapitulação}
% =====================================================

\begin{frame}{Onde estamos}
Duas tentativas de explicar a ciência --- ambas com problemas:
\vspace{0.4cm}
\pause
\begin{columns}[t]
\column{0.48\textwidth}
\begin{block}{Indutivismo}
Ciência \impt{começa pela observação}. \\[0.15cm]
Generaliza por \impt{indução}. \\[0.15cm]
Produz leis \impt{justificadas} pela experiência.
\end{block}

\column{0.48\textwidth}
\pause
\begin{block}{Falsificacionismo}
Ciência \impt{conjectura ousadamente}. \\[0.15cm]
Submete a \impt{testes severos}. \\[0.15cm]
Refuta --- nunca prova.
\end{block}
\end{columns}
\end{frame}

\begin{frame}{Indutivismo: os dois golpes}
\begin{itemize}
  \item \impt{Problema da indução}
  \begin{itemize}
    \item Sem justificação lógica para passar de singular a universal
    \item Justificar indução pela indução --- circular
  \end{itemize}
  \pause
  \vspace{0.3cm}
  \item \impt{Observação não é neutra}
  \begin{itemize}
    \item Carregada de teoria
    \item Não há ``afirmações singulares'' pré-teóricas
  \end{itemize}
\end{itemize}
\vspace{0.4cm}
\pause
\ntc{O que sobra da imagem indutivista?} Pouco.
\end{frame}

\begin{frame}{Falsificacionismo: a saída de Popper}
A \impt{assimetria lógica}:
\vspace{0.3cm}
\[
(T \rightarrow O) \wedge \neg O \;\vdash\; \neg T \qquad \text{(modus tollens)}
\]
\pause
\vspace{0.4cm}
\begin{itemize}
  \item Falsificação tem estatuto lógico privilegiado
  \item Critério de demarcação: \impt{falseabilidade}
  \item Versão sofisticada (Lakatos): predições novas corroboradas
\end{itemize}
\end{frame}

\begin{frame}{Mas o falsificacionismo também falha}
Três limitações graves (Cap.~VI):
\vspace{0.3cm}
\pause
\begin{enumerate}
  \item \impt{Observação carregada de teoria} \\
  Enunciados observacionais são \emph{falíveis}
  \pause
  \vspace{0.2cm}
  \item \impt{Tese de Duhem--Quine} \\
  $T \wedge H_1 \wedge \cdots \wedge H_n$ --- a falsificação não é conclusiva
  \pause
  \vspace{0.2cm}
  \item \impt{Inadequação histórica} \\
  Copérnico, atomismo, Newton e a Lua --- aparentes refutações \emph{ignoradas}
\end{enumerate}
\end{frame}

\begin{frame}{Teorias como Estruturas}
\begin{alertblock}{Mudança de perspectiva}
A unidade de análise da filosofia da ciência \\
\impt{não é a teoria isolada} --- é a \impt{estrutura teórica}.
\end{alertblock}
\vspace{0.4cm}
\pause
\begin{columns}[t]
\column{0.48\textwidth}
\begin{block}{Lakatos}
Programa de pesquisa \\[0.1cm]
\ntc{Lógico-metodológico}
\end{block}

\column{0.48\textwidth}
\pause
\begin{block}{Kuhn}
Paradigma \\[0.1cm]
\ntc{Histórico-sociológico}
\end{block}
\end{columns}
\end{frame}

% =====================================================
\section{Programas de pesquisa de Lakatos (Cap.~VII)}
% =====================================================

\begin{frame}{O ponto de partida}
Imre Lakatos: aluno e crítico de Popper.
\vspace{0.3cm}
\pause
\begin{itemize}
  \item Aceita: ciência avança por testes severos
  \pause
  \item Reconhece: cientistas \impt{não} abandonam teorias diante da primeira refutação
  \pause
  \item Conclusão: o objeto de avaliação não é a teoria isolada
\end{itemize}
\vspace{0.4cm}
\pause
\begin{block}{}
Avalia-se um \impt{programa de pesquisa} --- uma sequência articulada de teorias.
\end{block}
\end{frame}

\begin{frame}{Anatomia de um programa}
\begin{block}{Núcleo firme (\emph{hard core})}
Hipóteses fundamentais. \\
\ntc{Convencionalmente irrefutável} --- abandoná-lo é abandonar o programa.
\end{block}
\pause
\begin{block}{Cinturão protetor (\emph{protective belt})}
Hipóteses auxiliares, condições iniciais, teorias subsidiárias. \\
\ntc{Absorve} as refutações empíricas.
\end{block}
\pause
\begin{block}{Heurística}
\impt{Negativa:} ``não modifique o núcleo''. \\
\impt{Positiva:} ``modifique o cinturão neste e naquele sentido''.
\end{block}
\end{frame}

\begin{frame}{Exemplo: o programa newtoniano}
\begin{itemize}
  \item \impt{Núcleo:} as três leis do movimento $+$ lei da gravitação universal
  \pause
  \item \impt{Cinturão:} distribuição de massas, condições iniciais, atrito, instrumentos
  \pause
  \item \impt{Heurística positiva:} aplique a sistemas crescentemente complexos
  \begin{itemize}
    \item planeta isolado em torno do Sol
    \item perturbações entre planetas
    \item marés, rotação, \ldots
  \end{itemize}
\end{itemize}
\vspace{0.4cm}
\pause
\ntc{Anomalia em Urano?} Não se mexe em $F=ma$. \\
Postula-se um novo planeta: \impt{Netuno} (1846).
\end{frame}

\begin{frame}{Por que o núcleo é ``irrefutável''}
\begin{alertblock}{Não é defeito --- é decisão metodológica}
Sem núcleo estável, cada anomalia obriga a reconstruir tudo. \\
Com núcleo estável, há um problema bem definido: \\
\emph{como ajustar o cinturão para acomodar isto?}
\end{alertblock}
\pause
\vspace{0.4cm}
\ntc{Mas atenção:} se o cinturão tem de absorver demais, o programa começa a fazer ajustes \emph{ad hoc}. \\[0.2cm]
É aí que entra o critério de avaliação.
\end{frame}

\begin{frame}{Programas progressivos}
\begin{block}{Mudança progressiva (\emph{progressive problemshift})}
A sequência $T_1, T_2, T_3, \ldots$ é \impt{teoricamente progressiva} se cada $T_{i+1}$ \impt{prediz fatos novos}. \\[0.15cm]
É \impt{empiricamente progressiva} se algumas dessas predições são \impt{corroboradas}.
\end{block}
\pause
\vspace{0.4cm}
\ntc{Critério positivo:} o programa \emph{abre janelas} para fenômenos antes insuspeitos.
\end{frame}

\begin{frame}{Programas degenerativos}
\begin{block}{Mudança degenerativa (\emph{degenerative problemshift})}
Modificações sucessivas no cinturão apenas \impt{acomodam} fatos já conhecidos. \\[0.15cm]
Sem predições novas corroboradas. \\[0.15cm]
O programa \emph{fica para trás dos fatos}.
\end{block}
\pause
\vspace{0.4cm}
\ntc{Sintoma:} epiciclos sobre epiciclos. \\
Cada acréscimo salva um caso. Nada novo é predito.
\end{frame}

\begin{frame}{Newton vs.\ Ptolomeu tardio}
\begin{columns}[t]
\column{0.48\textwidth}
\begin{block}{Ptolomeu tardio}
Cada anomalia $\Rightarrow$ novo epiciclo \\[0.15cm]
\impt{Acomoda}, mas não prediz. \\[0.15cm]
\ntc{Degenerativo.}
\end{block}

\column{0.48\textwidth}
\pause
\begin{block}{Copérnico--Newton}
Achatamento da Terra \\
Retorno de cometas \\
Perturbações $\Rightarrow$ Netuno \\[0.15cm]
\ntc{Progressivo.}
\end{block}
\end{columns}
\vspace{0.5cm}
\pause
\impt{Critério lakatosiano:} um programa pode ser preferido a outro se for progressivo enquanto o rival é degenerativo.
\end{frame}

\begin{frame}{Comparação racional exige tempo}
A comparação entre programas rivais \impt{não é instantânea}.
\vspace{0.3cm}
\pause
\begin{itemize}
  \item Um programa em queda pode reerguer-se
  \pause
  \item Um programa em ascensão pode estagnar
  \pause
  \item É racional persistir em fase degenerativa --- por um tempo
\end{itemize}
\vspace{0.4cm}
\pause
\begin{alertblock}{}
Abandonar um programa é uma \impt{aposta científica} --- não uma dedução lógica.
\end{alertblock}
\end{frame}

\begin{frame}{O que Lakatos resolve}
\begin{itemize}
  \item Por que cientistas não abandonam teorias diante de anomalias \\
  \ntc{O cinturão é projetado para isto.}
  \pause
  \item Distingue modificação legítima de modificação \emph{ad hoc} \\
  \ntc{Pelo critério de progresso.}
  \pause
  \item Acomoda episódios como o copernicanismo nascente \\
  \ntc{Era racional perseguir o programa.}
  \pause
  \item Mantém critério \impt{objetivo} de progresso \\
  \ntc{Predições novas corroboradas.}
\end{itemize}
\end{frame}

\begin{frame}{O que Lakatos deixa em aberto}
\begin{itemize}
  \item \emph{Quando} um programa é claramente degenerativo? \\
  \ntc{Nunca há momento decisivo.}
  \pause
  \vspace{0.2cm}
  \item Como decidir entre programas que oscilam? \\
  \ntc{Pode levar décadas.}
  \pause
  \vspace{0.2cm}
  \item O modelo é \impt{prescritivo}, não tanto \impt{descritivo} \\
  \ntc{Como \emph{deve} ser avaliada a ciência --- não como ela funciona.}
\end{itemize}
\vspace{0.4cm}
\pause
A descrição empírica da prática científica foi proposta de modo mais radical por \impt{Kuhn}.
\end{frame}

% =====================================================
\section{Paradigmas de Kuhn (Cap.~VIII)}
% =====================================================

\begin{frame}{O olhar do historiador}
Thomas Kuhn é antes de tudo \impt{historiador da ciência}.
\vspace{0.3cm}
\pause
\begin{itemize}
  \item Estuda Copérnico, Lavoisier, Einstein
  \pause
  \item O que vê \impt{não se parece} com Popper
  \pause
  \item Cientistas \impt{não} passam o tempo tentando refutar suas teorias
\end{itemize}
\vspace{0.4cm}
\pause
\begin{block}{}
A maior parte do tempo, cientistas \impt{trabalham dentro} de uma estrutura aceita, resolvendo problemas técnicos.
\end{block}
\end{frame}

\begin{frame}{Os dois sentidos de paradigma}
\begin{block}{Sentido amplo --- matriz disciplinar}
Comprometimentos compartilhados por uma comunidade:
\begin{itemize}
  \item \impt{Generalizações simbólicas} --- $F = ma$
  \item \impt{Modelos e analogias} --- átomo como sistema solar
  \item \impt{Valores} --- precisão, simplicidade, fecundidade
  \item \impt{Exemplares} --- soluções-modelo
\end{itemize}
\end{block}
\pause
\begin{block}{Sentido estrito --- exemplar}
Solução-modelo de um problema. \\
Aprendendo a resolver \impt{exemplares}, o estudante \emph{internaliza} um paradigma.
\end{block}
\end{frame}

\begin{frame}{Paradigma vs.\ programa de pesquisa}
\begin{columns}[t]
\column{0.48\textwidth}
\begin{block}{Lakatos: programa}
Núcleo + cinturão + heurística \\[0.2cm]
\ntc{Ênfase:} lógica e metodologia \\
\ntc{Pergunta:} é progressivo?
\end{block}

\column{0.48\textwidth}
\pause
\begin{block}{Kuhn: paradigma}
Matriz disciplinar + exemplares \\[0.2cm]
\ntc{Ênfase:} comunidade e prática \\
\ntc{Pergunta:} resolve quebra-cabeças?
\end{block}
\end{columns}
\vspace{0.5cm}
\pause
\impt{Componente tácito} (treinamento, exemplares, valores compartilhados) é central em Kuhn.
\end{frame}

\begin{frame}{O ciclo da ciência}
\vspace{0.3cm}
\begin{center}
\begin{tikzpicture}[
  node distance=0.5cm,
  every node/.style={draw, rounded corners, fill=mprimary!10, draw=mprimary, minimum height=0.9cm, font=\small\bfseries, align=center},
  arrow/.style={-{Stealth[length=2.5mm]}, thick, draw=maccent}
]
\node (pre) {pré-ciência};
\node[right=of pre] (cn1) {ciência\\normal};
\node[right=of cn1] (crise) {crise};
\node[right=of crise] (rev) {revolução};
\node[right=of rev] (cn2) {nova ciência\\normal};

\draw[arrow] (pre) -- (cn1);
\draw[arrow] (cn1) -- (crise);
\draw[arrow] (crise) -- (rev);
\draw[arrow] (rev) -- (cn2);
\draw[arrow, dashed] (cn2.south) to[out=-90, in=-90] (crise.south);
\end{tikzpicture}
\end{center}
\vspace{0.3cm}
\pause
\ntc{O ciclo se repete.} Cada paradigma novo abre uma fase de ciência normal --- até a próxima crise.
\end{frame}

\begin{frame}{Pré-ciência}
\begin{itemize}
  \item Múltiplas escolas competindo
  \item Sem acordo sobre problemas legítimos
  \item Sem acordo sobre métodos
  \item Cada geração começa do zero
\end{itemize}
\vspace{0.4cm}
\pause
\begin{block}{Exemplo: ótica antes de Newton}
Aristotélicos, cartesianos, Hooke, Huygens --- pressupostos incompatíveis sobre a luz. \\
Após Newton: convergência (por algumas gerações) ao paradigma corpuscular.
\end{block}
\end{frame}

\begin{frame}{Ciência normal}
Estabelecido um paradigma: \impt{resolução de quebra-cabeças}.
\vspace{0.3cm}
\pause
\begin{itemize}
  \item Determinação precisa de fatos do paradigma
  \pause
  \item Comparação detalhada entre fato e teoria
  \pause
  \item Articulação --- novos casos, refinamentos, formulações elegantes
\end{itemize}
\vspace{0.4cm}
\pause
\begin{alertblock}{Crucial}
O paradigma \impt{não está sob teste}. \\
Sob teste está a \emph{habilidade do cientista} de fazer a natureza encaixar nele.
\end{alertblock}
\end{frame}

\begin{frame}{Anomalias}
Alguns quebra-cabeças resistem. Tornam-se \impt{anomalias}.
\vspace{0.3cm}
\pause
\begin{itemize}
  \item Inicialmente, \ntc{toleradas} --- nenhum paradigma resolve tudo
  \pause
  \item Tornam-se ameaça quando:
  \begin{itemize}
    \item atacam pressupostos centrais
    \item resistem a esforços prolongados
    \item bloqueiam aplicações importantes
  \end{itemize}
\end{itemize}
\vspace{0.4cm}
\pause
\begin{block}{}
Acumulando anomalias graves: \impt{crise}.
\end{block}
\end{frame}

\begin{frame}{Crise}
\begin{itemize}
  \item Confiança no paradigma se enfraquece
  \pause
  \item Surgem versões alternativas
  \pause
  \item Propostas heterodoxas começam a ser ouvidas
  \pause
  \item A fronteira do ``problema legítimo'' se borra
\end{itemize}
\vspace{0.4cm}
\pause
\begin{block}{Exemplo: física no fim do século XIX}
Michelson--Morley. Catástrofe do ultravioleta. Estabilidade dos átomos. \\[0.15cm]
Cada anomalia resistiu. O conjunto preparou \impt{relatividade} e \impt{quântica}.
\end{block}
\end{frame}

\begin{frame}{Revolução científica}
A crise é resolvida pela \impt{substituição} do paradigma.
\vspace{0.3cm}
\pause
\begin{itemize}
  \item Não é ajuste contínuo --- é \impt{descontinuidade}
  \pause
  \item Conceitos básicos \impt{mudam de sentido} (massa, espaço, tempo)
  \pause
  \item O que era problema legítimo deixa de ser
  \pause
  \item Novos exemplares e técnicas passam a definir o trabalho aceitável
\end{itemize}
\vspace{0.4cm}
\pause
\ntc{Depois:} nova ciência normal. O ciclo recomeça.
\end{frame}

\begin{frame}{Incomensurabilidade}
\begin{block}{Tese da incomensurabilidade}
Paradigmas rivais não podem ser comparados ponto a ponto por critérios neutros.
\end{block}
\pause
\vspace{0.4cm}
Três dimensões:
\begin{enumerate}
  \pause
  \item \impt{Semântica:} termos centrais mudam de sentido
  \pause
  \item \impt{Metodológica:} mudam os critérios do que conta como problema, método, explicação
  \pause
  \item \impt{Perceptual:} cientistas ``\emph{veem}'' o mundo de modos diferentes
\end{enumerate}
\end{frame}

\begin{frame}{Massa newtoniana e massa relativística}
\begin{columns}[t]
\column{0.48\textwidth}
\begin{block}{Newton}
Massa --- propriedade intrínseca \\
Conservada \\
Independe da velocidade
\end{block}

\column{0.48\textwidth}
\pause
\begin{block}{Einstein}
``Massa'' varia com velocidade \\
Massa de repouso é caso particular \\
$E = mc^2$
\end{block}
\end{columns}
\vspace{0.5cm}
\pause
Newton é \emph{aproximação} de Einstein para baixas velocidades.
\vspace{0.2cm}
\pause
\begin{alertblock}{}
Mas \impt{os conceitos não se traduzem perfeitamente}.
\end{alertblock}
\end{frame}

\begin{frame}{Atenção: incomensurável $\neq$ incomparável}
\begin{itemize}
  \item Cientistas de paradigmas rivais \impt{podem se entender}
  \pause
  \item Podem debater
  \pause
  \item Mas: não há \impt{linguagem neutra} a partir da qual avaliar os dois com critérios decisivos
\end{itemize}
\vspace{0.4cm}
\pause
\ntc{Não é caos.} É ausência de algoritmo de decisão.
\end{frame}

\begin{frame}{Como se escolhe entre paradigmas?}
A resposta de Kuhn é deliberadamente \impt{não algorítmica}:
\vspace{0.3cm}
\pause
\begin{itemize}
  \item Não há experimento decisivo
  \pause
  \item Não há critério lógico universal
  \pause
  \item Valores compartilhados (precisão, simplicidade, fecundidade) \\
  \ntc{podem ser ponderados de modo diferente}
  \pause
  \item Fatores \impt{extracientíficos} entram --- estética, intuição, contexto
\end{itemize}
\vspace{0.4cm}
\pause
\begin{block}{}
A escolha tem algo de \impt{conversão} --- mudança de gestalt sobre o domínio inteiro.
\end{block}
\end{frame}

\begin{frame}{Por que jovens lideram revoluções}
\begin{alertblock}{}
Revoluções são frequentemente lideradas por \impt{cientistas jovens} \\
ou por estranhos ao campo.
\end{alertblock}
\vspace{0.4cm}
\pause
\ntc{Razão:} estão menos comprometidos com o paradigma vigente. \\[0.3cm]
Não internalizaram décadas de exemplares.
\end{frame}

\begin{frame}{Kuhn é relativista?}
Lido por muitos como tal. Ele resistiu.
\vspace{0.3cm}
\pause
\begin{columns}[t]
\column{0.48\textwidth}
\begin{block}{Contra o relativismo}
Há \impt{progresso} dentro da ciência normal \\[0.1cm]
Valores compartilhados \\[0.1cm]
Revoluções não são arbitrárias
\end{block}

\column{0.48\textwidth}
\pause
\begin{block}{A favor do relativismo}
Sem critérios neutros decisivos \\[0.1cm]
Incomensurabilidade \\[0.1cm]
Fatores comunitários e estéticos
\end{block}
\end{columns}
\vspace{0.5cm}
\pause
\ntc{Leitura adequada:} intermediária. Rompe com a imagem da ciência convergindo cumulativamente para a verdade --- mas \emph{não} é relativismo radical.
\end{frame}

\begin{frame}{O que Kuhn resolve}
\begin{itemize}
  \item A \impt{ciência normal} é a atividade científica predominante
  \pause
  \item Revoluções são \impt{descontinuidades estruturais}, não acumulação suave
  \pause
  \item Reabilita a \impt{história} da ciência como recurso filosófico
\end{itemize}
\end{frame}

\begin{frame}{O que Kuhn deixa em aberto}
\begin{itemize}
  \item Se paradigmas são incomensuráveis, em que sentido há \impt{progresso}?
  \pause
  \vspace{0.2cm}
  \item Como distinguir escolha racional de mera \impt{conversão} ideológica?
  \pause
  \vspace{0.2cm}
  \item A categoria ``paradigma'' é clara o bastante? \\
  \ntc{Críticos identificaram mais de 20 sentidos.}
  \pause
  \vspace{0.2cm}
  \item Funciona bem para física madura --- e para ciências sociais? Computação?
\end{itemize}
\end{frame}

% =====================================================
\section{Síntese}
% =====================================================

\begin{frame}{Lakatos vs.\ Kuhn}
\begin{columns}[t]
\column{0.48\textwidth}
\begin{block}{Lakatos --- Programa}
Núcleo + cinturão + heurística \\[0.15cm]
\impt{Lógico-metodológico} \\[0.15cm]
Critério: progressivo? \\[0.15cm]
\ntc{Continuidade} entre programas
\end{block}

\column{0.48\textwidth}
\begin{block}{Kuhn --- Paradigma}
Matriz disciplinar + exemplares \\[0.15cm]
\impt{Histórico-sociológico} \\[0.15cm]
Critério: resolve quebra-cabeças? \\[0.15cm]
\ntc{Descontinuidade} revolucionária
\end{block}
\end{columns}
\vspace{0.5cm}
\pause
\impt{Diagnóstico comum:} a unidade de análise é a estrutura, não a teoria isolada.
\end{frame}

\begin{frame}{Tensão central}
A unidade de análise é a mesma --- a estrutura, não a teoria isolada. \\
A pergunta filosófica é: \impt{a ciência ainda é racional e objetiva?}
\vspace{0.1cm}
\pause
\begin{columns}[t]
\column{0.48\textwidth}
\begin{block}{Lakatos diz: sim}
A racionalidade está na \impt{história} do programa. \\[0.2cm]
\ntc{Preserva} a objetividade --- só relocaliza-a no tempo. \\[0.2cm]
Critério: predições novas corroboradas.
\end{block}

\column{0.48\textwidth}
\pause
\begin{block}{Kuhn diz: depende}
A escolha entre paradigmas envolve \impt{julgamento comunitário}. \\[0.2cm]
\ntc{Compromete} a imagem de uma decisão puramente lógica. \\[0.2cm]
Sem algoritmo de decisão.
\end{block}
\end{columns}
\vspace{0.5cm}
\pause
\begin{alertblock}{}
\centering
A pergunta de fundo: \impt{a ciência converge à verdade}, ou apenas \impt{resolve mais problemas}?
\end{alertblock}
\end{frame}
\begin{frame}{Lição geral}
Ambos forçam a filosofia da ciência a ir \impt{além da lógica}.
\vspace{0.4cm}
\pause
Onde o falsificacionismo via \emph{teoria $\times$ experimento}:
\vspace{0.2cm}
\begin{itemize}
  \item Lakatos vê \impt{programas}
  \item Kuhn vê \impt{paradigmas}
\end{itemize}
\pause
\ntc{Considerar:} história, comunidade, prática, formação. \\[0.3cm]
A ciência é uma atividade humana --- e qualquer filosofia adequada precisa dar conta disso.
\end{frame}

\begin{frame}[allowframebreaks]{Referências}
\begin{thebibliography}{9}
\bibitem[Chalmers(1993)]{chalmers1993}
Chalmers, A.\ F. (1993).
\newblock \emph{O que é ciência afinal?} (2ª ed.).
\newblock São Paulo: Brasiliense.

\bibitem[Kuhn(1962)]{kuhn1962}
Kuhn, T.\ S. (1962).
\newblock \emph{The structure of scientific revolutions}.
\newblock Chicago: University of Chicago Press.

\bibitem[Lakatos(1978)]{lakatos1978}
Lakatos, I. (1978).
\newblock \emph{The methodology of scientific research programmes}.
\newblock Cambridge: Cambridge University Press.

\bibitem[Popper(1972)]{popper1972}
Popper, K.\ R. (1972).
\newblock \emph{A lógica da pesquisa científica}.
\newblock São Paulo: Cultrix.
\end{thebibliography}
\end{frame}

\end{document}
